\section{The next turn should inherit a better means}
\label{sec:introduction}

Software maintenance repeatedly pays for the same kinds of work: locating committed inputs, reconstructing dependency information, preparing evidence, and explaining why a result is safe to use. An optimization can remove one such cost. A more consequential outcome is a qualified engineering instrument that helps another maintainer make a different improvement. The central question is whether this retained instrument improves fresh work after the costs of discovering, constructing, checking, integrating and maintaining it are included.

\begin{quote}
\textbf{Build software that improves useful software, and leaves behind better means of making the next improvement.}
\end{quote}

The Software Capability Foundry examines that question through \emph{qualified boundaries as reusable engineering assets}. A boundary separates the implementation that may change from the externally required behavior that must remain valid. Its value is practical: a later constructor can use a well-defined operation without reconstructing every implementation decision, while a verifier can determine when that operation is no longer applicable. The asset is therefore executable code together with a checked contract, provenance, resource account, failure behavior and a way back to the incumbent. A saved answer or a successful benchmark row does not provide this interface.

The immediate setting is repository and publication tooling in the public AGI ALPHA ecosystem. Repeated short-lived tool invocations are a plausible source of overhead, but invocation count alone is not a performance result. The study begins with pinned source call sites and reproduces their required outputs in disposable local repositories. It distinguishes useful restructuring of the computation from a relaxation of checks. For example, changing how committed bytes are acquired must preserve which revision, object type and content the caller receives; skipping a release-integrity check would change the mission.

\subsection{Three achievements require three tests}

We distinguish a \emph{service improvement}, a \emph{transferable capability}, and an \emph{improved improvement process}. The first delivers a useful result at lower justified burden. The second helps on a fresh problem with changed inputs or structure. The third reduces the complete work needed to discover or qualify a further improvement. These implications do not follow automatically. Reusing a fast function may establish the second without establishing the third. A child program that imports a parent without benefiting from it establishes neither.

The visible journey contains four predeclared missions. Mission A delivers a useful improvement and retains its engineering interface. Mission B changes the workload and tests that interface's contribution. Mission C uses accepted prior work to help construct or qualify a distinct capability with a different output contract. Mission D changes the operating conditions enough that reuse must be refused, replaced or requalified. These explanatory missions remain separate from the protected confirmation instances.

The decisive comparisons concern useful content rather than lineage terminology. Removing the parent still permits reconstruction from legitimate knowledge and ordinary repository history. An independently constructed equivalent asset tests whether any benefit belongs to the asset's content. A mature specialist and an equally enabled conventional tool builder remain eligible to win. This follows the baseline and evidence-package discipline of AGI ALPHA rather than treating a new label as a privileged comparator~\cite[PDF pp.~37, 60, 69]{agi_alpha}.

\subsection{Why the earlier failures matter}

The preceding research supplies useful mechanisms and reasons for stricter comparisons. The Possibility Foundry constructed a reusable instrument and used it in a distinct temporal capability, but an analytical specialist and an independent construction policy found the same instrument more cheaply. Charging the complete ancestor search to the small successor campaign also reversed its affordability~\cite[PDF pp.~2--4]{possibility_foundry}.

The published coordination record for Strategy Succession reports 517 qualified missions out of 864 for the primary flywheel, tied with the equally enabled learner and below the stronger portfolio's 591. No primary retained capability was admitted; no fresh repair output both invoked its parent and qualified. Its diagnostic inheritance advantage against exact reconstruction disappeared against a cheaper conventional fallback~\cite{strategy_record}. The subsequent \emph{Earn the Next Turn} record reports constructive succession and parent-invoking fresh outputs, while its whole-participant CPU comparison failed the predeclared five-percent admission margin against native factor reuse and refactorization~\cite{earned_turn_record}. These are preserved prior results, not observations from the new experiment.

Those findings motivate a narrower and more demanding contribution. The present study asks an engineering asset to survive a competent alternative, a current correctness check, complete participant accounting and an explicit changed-scope case. It does not force parent reuse when a cheaper direct method works. Service resources and improvement resources are recorded separately, so unsuccessful research does not silently consume the resources promised for the useful delivery.

\subsection{Scope of the contribution}

The experiment studies bounded deterministic construction and qualification. Its primitive tools and candidate mechanisms are supplied and disclosed. Any optimization found within that space is a result about the executed search and selected artifact; it is not evidence of general autonomous invention. The browser experience makes the goal, result, retained interface, comparison and decision inspectable, while distinguishing live local computation, recorded experiment evidence and execution through an approved external adapter.

The work is additive to the existing workbench and mission history. It introduces no production scheduler and derives no execution authority from a research result. The empirical report must answer two separate questions: whether the capability is worth using on the next specified mission, and whether its complete investment has paid for itself. Unknown monetary or human costs remain unknown. A reliable reduction in engineering burden can be useful without demonstrating cash profit, open-ended compounding or an increasing rate of improvement.
